\documentclass[11pt]{article}
\usepackage[margin=0.86in]{geometry}
\usepackage[T1]{fontenc}
\usepackage{lmodern}
\usepackage{amsmath,amssymb,amsthm,graphicx,microtype}
\usepackage[hidelinks]{hyperref}

\newtheorem{theorem}{Theorem}
\newtheorem{lemma}[theorem]{Lemma}
\newcommand{\R}{\mathbb R}
\newcommand{\N}{\mathbb N}
\newcommand{\cP}{\mathcal P}

\title{Strict Positivity on Even Powers Does Not Imply\\
Square Positivity Without Continuity}
\author{Candidate counterexample packet\\
\small arXiv:1610.01659, Question 2.9}
\date{11 August 2026}

\begin{document}
\maketitle

\noindent\textbf{Status.} Candidate full negative answer to the
continuity-\emph{removal} part of Question 2.9, for every fixed $d\ge2$,
subject to expert review.  The construction does not classify the additional
hypotheses weaker than continuity under which the implication might be
recovered.

\section{Source question}

Infusino and Kuhlmann prove that if $(A,\tau)$ is a locally multiplicatively
convex real algebra and $L:A\to\R$ is $\tau$-continuous, then the following
are equivalent: $L$ is nonnegative on sums of squares; $L$ is nonnegative on
sums of $2d$-th powers for every $d$; and it is nonnegative on such powers for
at least one $d$ \cite[Corollary 2.8]{IK}.  They then ask on printed page 6:
\begin{quote}
Does Corollary 2.8 still hold if we remove or weaken the continuity
assumption?
\end{quote}

\begin{center}
\includegraphics[width=0.94\textwidth]{figures/open_problem_crop.png}
\end{center}

We show that continuity cannot simply be removed.  The counterexample already
occurs in the one-variable polynomial algebra.

\section{Proof intuition}

A functional on $\R[x]$ is determined by its moments $m_j=L(x^j)$.  Positivity
on $2d$-th powers says that every finite Hankel form
\[
 (c_0,\ldots,c_N)\longmapsto
 L\!\left(\left(\sum_{j=0}^N c_jx^j\right)^{2d}\right)
\]
is positive semidefinite.  This higher-order condition need not make the
ordinary Hankel matrix positive semidefinite.

We start in dimension two with a strictly positive binary form whose moment
$m_2$ is negative.  To add one polynomial degree, first declare all new
moments except the very last one to be zero.  The extended form is the old
positive-definite form plus finitely many mixed terms plus
$M|c_{N+1}|^{2d}$.  On the old finite-dimensional coefficient space, the old
form has a uniform positive lower bound and each mixed form has a uniform
upper bound.  Young's inequality absorbs every mixed term into those two pure
terms once the new top moment $M$ is sufficiently large.  Iteration gives one
functional on all polynomials.

\section{A positive-definite Hankel-form completion lemma}

Fix $d\ge2$ and write $r:=2d$.  For real numbers $m_0,m_1,\ldots$, define
\[
 P_N(c_0,\ldots,c_N)
 :=\sum_{i_1,\ldots,i_r=0}^N
 m_{i_1+\cdots+i_r}c_{i_1}\cdots c_{i_r}.
 \tag{1}
\]

\begin{lemma}\label{lem:completion}
There is a real sequence $(m_j)_{j\ge0}$ such that every $P_N$ is positive
definite, $m_2=-\binom{r}{2}^{-1}<0$, and the subsequence $(m_{rN})_N$ is
unbounded.
\end{lemma}

\begin{proof}
For $0\le j\le r$, set
\[
 m_0=m_r=1,\qquad m_2=-\binom r2^{-1},\qquad
 m_j=0\quad(j\notin\{0,2,r\}).
\]
Then
\[
 P_1(a,b)=a^r-a^{r-2}b^2+b^r.
\]
Weighted AM--GM gives
\[
 |a|^{r-2}|b|^2
 \le \frac{d-1}{d}|a|^r+\frac1d|b|^r,
\]
and therefore
\[
 P_1(a,b)\ge \frac1d|a|^r+\frac{d-1}{d}|b|^r.
 \tag{2}
\]
Thus $P_1$ is positive definite.

Suppose moments through $m_{rN}$ have been chosen and $P_N$ is positive
definite.  Define
\[
 m_{rN+1}=\cdots=m_{r(N+1)-1}=0,
 \qquad m_{r(N+1)}=M,
 \tag{3}
\]
where $M>0$ will be fixed below.  If $q\in\R[x]$ has degree at most $N$, the
binomial expansion has the form
\begin{align}
 L\bigl((q+t x^{N+1})^r\bigr)
 =P_N(q)+\sum_{k=1}^{r-1}t^k A_{k,N}(q)+Mt^r.       \tag{4}
\end{align}
Here $A_{k,N}$ is already determined before $M$ is chosen: its largest
possible moment index is
\[
 k(N+1)+(r-k)N=rN+k<r(N+1).
\]
It is a homogeneous polynomial of degree $r-k$ in the coefficients of $q$.

Give the space of polynomials of degree at most $N$ any Euclidean coefficient
norm.  Compactness of its unit sphere and positive definiteness give
\[
 P_N(q)\ge\delta_N\|q\|^r
\]
for some $\delta_N>0$.  Likewise, there are constants $C_{k,N}$ with
$|A_{k,N}(q)|\le C_{k,N}\|q\|^{r-k}$.  The weighted Young inequality provides
numbers $D_{k,N}<\infty$ such that, for all $q$ and $t$,
\[
 C_{k,N}|t|^k\|q\|^{r-k}
 \le \frac{\delta_N}{2(r-1)}\|q\|^r+D_{k,N}|t|^r.
 \tag{5}
\]
Choose
\[
 M>\sum_{k=1}^{r-1}D_{k,N}+N+1.
\]
Equations (4)--(5) yield
\[
 P_{N+1}(q+t x^{N+1})
 \ge \frac{\delta_N}{2}\|q\|^r
   +\left(M-\sum_{k=1}^{r-1}D_{k,N}\right)|t|^r>0
\]
unless $q=t=0$.  This completes the induction.  The additional $N+1$ in the
choice of $M$ makes $(m_{rN})_N$ unbounded.
\end{proof}

\section{Main counterexample}

\begin{theorem}\label{thm:main}
For every integer $d\ge2$ there exist a unital commutative normed real algebra
$(A,\rho)$ and a linear functional $L:A\to\R$ such that
\[
 L(q^{2d})>0\quad\hbox{for every }q\ne0,
 \qquad L(x^2)<0.
 \tag{6}
\]
Consequently, condition (3) of Corollary 2.8 in \cite{IK} does not imply
condition (1) after continuity is removed, and $L$ has no nonnegative
representing measure.
\end{theorem}

\begin{proof}
Take $A=\R[x]$ with the submultiplicative coefficient norm
\[
 \rho\!\left(\sum_ja_jx^j\right):=\sum_j|a_j|.
\]
Let $(m_j)$ be supplied by Lemma~\ref{lem:completion} and define the linear
functional $L$ by $L(x^j)=m_j$.  Every nonzero polynomial $q$ belongs to a
finite degree space, so (1) and positive definiteness of the corresponding
$P_N$ give $L(q^{2d})>0$.  Hence $L$ is nonnegative on
$\sum A^{2d}$.  On the other hand,
\[
 L(x^2)=m_2=-\binom{2d}{2}^{-1}<0,
\]
so $L$ is not nonnegative on $\sum A^2$.

The functional is not $\rho$-continuous: continuity would imply
$|m_j|=|L(x^j)|\le\|L\|\rho(x^j)=\|L\|$ for all $j$, contradicting the
unbounded subsequence in Lemma~\ref{lem:completion}.  Finally, a nonnegative
representing measure would give $L(x^2)=\int x^2\,d\mu\ge0$, which is
impossible.
\end{proof}

\section{Verification, scope, and novelty}

The proof was checked at three levels.  First, (2) is an explicit coercive
bound for the negative-second-moment base case.  Second, the degree identity
$rN+k<r(N+1)$ guarantees that the coefficient $M$ occurs only in the pure
new-coordinate term.  Third, the finite-dimensional compactness and Young
bounds in (5) make the induction uniform over all old polynomials, not merely
pointwise.  The accompanying script verifies the base coefficients and this
degree bookkeeping for $2\le d\le10$ and the first thirty extension stages.

This settles the version with continuity deleted, for every $d\ge2$.  It does
not decide which particular quasi-analytic, boundedness, or other hypotheses
strictly weaker than continuity may restore the implication.

A bounded novelty search through 11 August 2026 covered the run indexes, the
exact text of Question 2.9, the source title and authors, later infinite
moment-problem papers, and combinations of ``positive on $2d$-th powers,''
``square positive,'' ``Hankel form,'' and ``continuity.''  The search found
finite-dimensional positive semidefinite Hankel tensors whose associated
Hankel matrices are not positive semidefinite \cite{WLQX,DQW}, but no global
linear functional on $\R[x]$, recursive all-degree completion, or stated
answer to Question 2.9.  This is a bounded retrieval result, not proof of
novelty.

\medskip
\noindent\textbf{Human-review recommendation.} High priority.  Check the
uniform absorption in (5) and the moment-index calculation in (4).  Those are
the only nontrivial completion steps; once established, (6) is immediate.

\begin{thebibliography}{9}
\bibitem{IK}
M. Infusino and S. Kuhlmann,
\emph{Infinite dimensional moment problem: open questions and applications},
Contemp. Math. \textbf{697} (2017), 187--201; arXiv:1610.01659.

\bibitem{WLQX}
Q. Wang, G. Li, L. Qi, and Y. Xu,
\emph{New classes of positive semi-definite Hankel tensors},
Minimax Theory Appl. \textbf{2} (2017), 231--248; arXiv:1411.2368.

\bibitem{DQW}
W. Ding, L. Qi, and Y. Wei,
\emph{Inheritance properties and sum-of-squares decomposition of Hankel
tensors: theory and algorithms}, arXiv:1505.02528.
\end{thebibliography}

\end{document}
